\documentclass[10pt,a4paper]{article}
\usepackage[a4paper,left=1.8cm,right=1.8cm,top=2.0cm,bottom=1.8cm]{geometry}
\usepackage[T1]{fontenc}
\usepackage[scaled=0.94]{helvet}
\renewcommand{\familydefault}{\sfdefault}
\usepackage{amsmath,amssymb,bm}
\usepackage{xcolor}
\usepackage{booktabs}
\usepackage{fancyhdr}
\usepackage{titlesec}
\usepackage{enumitem}
\usepackage[numbers,sort&compress]{natbib}
\definecolor{ink}{HTML}{1B2A41}
\definecolor{accent}{HTML}{1F5F8B}
\usepackage[colorlinks=true,linkcolor=accent,citecolor=accent,urlcolor=accent]{hyperref}
\usepackage{url}

\titleformat{\section}{\Large\bfseries\color{ink}}{\thesection.}{0.6em}{}
\titleformat{\subsection}{\large\bfseries\color{accent}}{\thesubsection}{0.6em}{}
\titleformat{\subsubsection}{\normalsize\bfseries\color{ink}}{\thesubsubsection}{0.6em}{}
\titlespacing*{\section}{0pt}{1.4ex plus .5ex}{0.8ex}
\titlespacing*{\subsection}{0pt}{1.2ex plus .4ex}{0.5ex}

\pagestyle{fancy}
\fancyhf{}
\fancyhead[L]{\small\color{accent}xphd handbook --- Part II: methods and derivations}
\fancyfoot[R]{\small\color{gray}\thepage}
\renewcommand{\headrulewidth}{0.4pt}
\renewcommand{\headrule}{\hbox to\headwidth{\color{gray!40}\leaders\hrule height \headrulewidth\hfill}}

\setlength{\parskip}{0.45ex}
\setlength{\parindent}{0pt}
\newcommand{\tbd}[1]{\textcolor{red}{[#1]}}
\newcommand{\G}{\mathcal{G}}
\newcommand{\bq}{\mathbf{q}}
\newcommand{\bQ}{\mathbf{Q}}
\newcommand{\bk}{\mathbf{k}}
\newcommand{\bK}{\mathbf{K}}
\newcommand{\bb}{|\mathbf{b}|}

\begin{document}
\setcounter{page}{25}
\setcounter{section}{11}

\thispagestyle{fancy}
{\color{ink}\huge\bfseries Part II --- Methods and derivations\par}
\vspace{0.6ex}
{\color{accent}\large The physics and numerics behind the xphd results: every
formula the code evaluates, how it is derived, and how each was checked\par}
\vspace{1.2ex}

Part~I (pages 1--24) explains the package: the commands, scripts, data files and
material walkthroughs. Part~II collects the methods behind them. Each section
states the quantity, derives the working formula, describes how it is evaluated
on the mesh, and gives the checks that established it, with the numbers
obtained for GaN, hBN and WSe$_2$. Results that are still pending are marked
\tbd{in red}. The references at the end carry their verification status in the
accompanying \texttt{partII.bib}: entries marked VERIFY there should be checked
before they are cited in a paper.

\vspace{0.6ex}
{\setlength{\parskip}{0pt}\small\tableofcontents}

\section*{Errata to Part I (4 October 2026)}
\addcontentsline{toc}{section}{Errata to Part I}
\begin{enumerate}[nosep,leftmargin=1.4em]
\item \textbf{No \texttt{KfnQPdb} fault.} Sections 6.4, 9.1, 9.2 and the pitfalls
table of Sec.~10 state that \texttt{Yambo}'s \texttt{KfnQPdb} applied the
quasiparticle corrections without the Kohn-Sham energies. That is wrong. With
\texttt{ExtendOut}, \texttt{o-*.qp} lists \texttt{K-point, Band, Eo, E, E-Eo,
Vxc, Vnlxc, Sc|Eo, \ldots}: the fourth column is the quasiparticle energy $E$,
and reading it as $E-E_o$ produced the false diagnosis. \texttt{KfnQPdb} applies
\texttt{ndb.QP} correctly; the report's \texttt{[K+QP]} gap is the true
quasiparticle gap.
\item \textbf{The numbers of Sec.~9.1.} The values quoted there as corrections
at $\bK$ ($-0.4709$, $+0.8867$, $+0.8716$~eV for bands 62, 63, 64) are the
quasiparticle energies. The corrections of that (RIM-W) GW are $-0.471$,
$-0.373$ and $-0.425$~eV: it opened the gap at $\bK$ by only $+0.098$~eV (the
report's 1.3576~eV) and put band 64 below band 63 by 15~meV (not ``+37.2 to
+22.2~meV''). That ordering, not a code fault, is why the BSE with the database
put the bright exciton lowest.
\item \textbf{The scissor.} The 1.357578~eV scissor of the WSe$_2$ BSE runs is that
GW's quasiparticle gap at $\bK$, not its correction; those runs used a gap of
2.617~eV at $\bK$. The GW without RIM-W gives a correction of $+1.181$~eV
(gap 2.440~eV for 62$\to$63, 2.422~eV for 62$\to$64).
\item \textbf{RIM-W, scope.} Sections 6.4 and 10 of Part~I advise dropping
\texttt{rim\_w} in general. That holds for WSe$_2$ only: there it broke the BSE
kernel and the GW (gap correction at $\bK$ $+0.10$~eV with it, $+1.18$~eV without),
whereas GaN and hBN ran correctly with RIM-W. The cause is not isolated (spinors,
the $12\times12$ mesh and the RIM-W settings all differed). Check the plausibility
of the gap correction and of the binding energy; if in doubt, compare a run
without \texttt{rim\_w}.
\item \textbf{Tools.} \texttt{cb\_splitting.py} read \texttt{o-*.qp} by position and
misread \texttt{ExtendOut} files; it now reads columns by header name, as
\texttt{qp\_decompose.py} does, and reports band edges and the direct or indirect
gap (Sec.~\ref{sec:gw}).
\end{enumerate}

% =====================================================================
\section{Notation and conventions}
% =====================================================================
\label{sec:notation}

\textbf{Lattice.} The hexagonal reciprocal basis $\mathbf{b}_1,\mathbf{b}_2$
encloses 60$^\circ$; reduced coordinates $(q_1,q_2)$ give
$\bq=q_1\mathbf{b}_1+q_2\mathbf{b}_2$, and lengths are quoted in units of
$|\mathbf{b}|=|\mathbf{b}_1|$. In these units the Brillouin zone has area
$A_{\mathrm{BZ}}=\sqrt3/2$, $\bK=(1/3,1/3)$, $\bK'=(2/3,2/3)$ and the
$\mathbf{M}$ points are the edge midpoints, e.g.\ $(1/2,0)$. The source mesh is
$N\times N$ ($N=24$ for GaN and hBN, 12 for WSe$_2$), with spacing $h=1/N$; the
refined mesh used for integration is $360\times360$.

\textbf{Electron-phonon elements.} Computed by \texttt{LetzElPhC}
\cite{LetzElPhC} in the \texttt{Yambo} convention,
$g_{n'n\nu}(\bk,\bq)=\langle\psi_{n'\bk}|\Delta_{\nu\bq}V|\psi_{n,\bk-\bq}\rangle$,
and converted to the standard convention before projection. The phonon
displacement normalisation is that of DFPT \cite{Baroni2001,Giustino2017}.

\textbf{Exciton envelopes.} The BSE is solved with \texttt{Yambo}
\cite{Sangalli2019} at every irreducible $\bQ$: for GaN with the
subspace-iteration solver of SLEPc \cite{Hernandez2005}, for WSe$_2$ by full
diagonalisation (matrices of dimension 2304). We write $A^{S\bQ}_{vc\bk}$ with the electron at
$\bk$ and the hole at $\bk-\bQ$. \texttt{Yambo}'s BSE tables store the electron
at the table index $\bk$ and the hole at $\bk+\bQ$, i.e.\ the exciton at $-\bQ$
in this convention. Time reversal gives the excitons at $\bQ$ and $-\bQ$ the same
energies and parities, so nothing physical changes, but every formula that
pairs an electron and a hole must use the stored pairing (Sec.~\ref{sec:excparity}).

\textbf{Archives.} An archive \texttt{GI\_ExcPh\_QXXXX.npz} stores, for one
exciton momentum $\bQ$, the complex amplitudes
$G[\bq,\nu,\lambda,\beta]=\langle\beta,\bQ+\bq|\partial V_\nu|\lambda,\bQ\rangle=\G_{\beta\lambda\nu}(\bQ,\bq)$
in eV, for every $\bq$ of the source mesh, every branch $\nu$, initial state
$\lambda$ at $\bQ$ and final state $\beta$ at $\bQ+\bq$, together with
$|\G|^2$, its electron and hole parts, and the energies. The final envelope is
the conjugated one, as in \texttt{Yambopy} \cite{Yambopy}.

% =====================================================================
\section{Exciton-phonon matrix elements}
% =====================================================================
\label{sec:excph}

\subsection{Definition}
The exciton-phonon matrix element is the sum of a hole and an electron term, in
each of which one carrier scatters while the other is a spectator
\cite{Antonius2022,Chen2020}:
\begin{equation}
\begin{split}
\G_{\beta\lambda\nu}\!\left(\bQ,\bq\right) ={}&
\sum_{v,v',c,c',\bk}
A_{\lambda,\bQ}^{v,c,\bk}
\!\left[g_{vv',\nu}\!\left(\bk-\bQ,\bq\right)\delta_{c,c'}\right]
\left(A_{\beta,\bQ+\bq}^{v',c',\bk}\right)^{\ast}\\
&-\sum_{v,v',c,c',\bk}
A_{\lambda,\bQ}^{v,c,\bk}
\!\left[g_{c'c,\nu}^{\ast}\!\left(\bk+\bq,\bq\right)\delta_{v,v'}\right]
\left(A_{\beta,\bQ+\bq}^{v',c',\bk+\bq}\right)^{\ast}.
\end{split}
\label{eq:G}
\end{equation}

\subsection{Gauge}
Both factors in Eq.~(\ref{eq:G}) are built from the same Kohn-Sham wavefunctions
(one \texttt{Yambo} SAVE for \texttt{LetzElPhC} and the BSE), so they share one
single-particle gauge. Two freedoms remain on the exciton side.

\emph{Overall phases.} Replacing $A^{\lambda}\to e^{i\theta_\lambda}A^\lambda$ and
$A^{\beta}\to e^{i\theta_\beta}A^\beta$ multiplies $\G_{\beta\lambda\nu}$ by
$e^{i(\theta_\lambda-\theta_\beta)}$, which cancels in $|\G_{\beta\lambda\nu}|^2$.

\emph{Degenerate manifolds.} Within a degenerate manifold $\mathcal{M}$ the
solver's basis is arbitrary: $|\beta'\rangle=\sum_{\beta\in\mathcal{M}}U_{\beta\beta'}|\beta\rangle$
with $U$ unitary. Then $\sum_{\beta'\in\mathcal{M}}|\G_{\beta'\lambda\nu}|^2=\sum_{\beta\in\mathcal{M}}|\G_{\beta\lambda\nu}|^2$,
and likewise for the initial manifold. Only sums over complete manifolds are
basis-independent, and every invariant used below is such a sum.

\subsection{Excitons on the full zone}
The BSE is solved only at the irreducible momenta (61 for a $24\times24$ mesh,
19 for $12\times12$). The exciton at any other momentum is the image of its
irreducible parent under the operation $R$ that maps one to the other
\cite{Nalabothula2026}:
\begin{equation}
A^{S,R\bQ}_{vc,R\bk}=\sum_{v'c'}D^{c}_{cc'}(R;\bk)\,A^{S\bQ}_{v'c',\bk}\,
\bigl[D^{v}_{vv'}(R;\bk-\bQ)\bigr]^{*},
\label{eq:rotate}
\end{equation}
where $D^{n}(R;\bk)$ are the single-particle rotation matrices in the BSE band
window, computed by \texttt{LetzElPhC}. For an operation that includes time
reversal the parent envelope is complex conjugated first. The point is not to
fix phases -- they cancel anyway -- but to make every member of a star the exact
symmetry image of the same parent; a mapping that does not do so breaks the
symmetry of $|\G|^2$.

\subsection{Symmetry tests}
\emph{Star uniformity.} The integrated strength
$T(\bq)=\sum_{\nu\lambda\beta}|\G_{\beta\lambda\nu}(\bQ=0,\bq)|^2$, summed over
all branches and complete manifolds, must be identical for every member of a
star \cite{Dresselhaus2008,BradleyCracknell1972,Maradudin1968}. For GaN the
ratio $T_{\max}/T_{\min}$ within a star is 1.0001 for the reconstructed full
zone, against 2.25 for an initial output whose excitons at the star members were
not constructed as images of a common parent.

\emph{Manifold-resolved test} (\texttt{xphd check-archive}). For an archive at a
time-reversal-invariant $\bQ$, define for each $\bq$
$S_{abc}(\bq)=\sum_{\nu\in a}\sum_{\lambda\in b}\sum_{\beta\in c}|\G_{\beta\lambda\nu}(\bQ,\bq)|^2$
over phonon manifolds $a$ (branches within 0.1~meV), initial manifolds $b$ and
final manifolds $c$ (states within a tolerance, default 15~meV). Symmetry
requires $S(\bq)=S(g\bq)$ for $g=$ time reversal ($\bq\to-\bq$) and $g=C_3$. The
reported deviation is
\begin{equation}
L_1=\frac{\sum_{\bq}\sum_{abc}|S_{abc}(\bq)-S_{abc}(g\bq)|}
{\tfrac12\sum_{\bq}\sum_{abc}\bigl[S_{abc}(\bq)+S_{abc}(g\bq)\bigr]} ,
\end{equation}
with OK below $10^{-2}$ and FAIL above $5\times10^{-2}$. The C$_3$ matrix in
reduced coordinates depends on the angle between $\mathbf{b}_1$ and
$\mathbf{b}_2$: $(q_1,q_2)\to(-q_1-q_2,q_1)$ for 60$^\circ$ and
$(-q_2,q_1-q_2)$ for 120$^\circ$; the code takes the one that leaves the exciton
energies invariant. Results: GaN $1.3\times10^{-3}$ (time reversal) and
$3.3\times10^{-3}$ (C$_3$); WSe$_2$ $1.2\times10^{-3}$ and $1.0\times10^{-3}$. A
rotation fault shows as one operation far above the other.

\emph{Exchange-split doublets.} At $\bQ=0$ the BSE evaluates the long-range
exchange along one direction $\hat{\mathbf{e}}$, splitting each E$'$ doublet
into members polarised parallel and perpendicular to $\hat{\mathbf{e}}$ (GaN:
2.8~meV for the first bright doublet, 10~meV for the second). For a member with
real in-plane dipole $\mathbf{d}$, the coupling to small-$\bq$ polar phonons
carries the factor $|\hat\bq\cdot\mathbf{d}|^2$. It is unchanged by
$\bq\to-\bq$ but not by C$_3$, which rotates $\hat\bq$ relative to the fixed
$\mathbf{d}$. Summed over the doublet, $|\hat\bq\cdot\mathbf{d}_\parallel|^2+
|\hat\bq\cdot\mathbf{d}_\perp|^2=|\mathbf{d}|^2$ is isotropic. With 5~meV
manifolds the GaN second doublet is split in two and the C$_3$ test fails
($L_1=0.20$, all of it in that doublet and the LO branch at the smallest
$|\bq|$); grouping it restores $3.3\times10^{-3}$.

% =====================================================================
\section{Mirror parity and selection rules}
% =====================================================================
\label{sec:parity}

\subsection{The horizontal mirror}
$\sigma_h$ maps $(x,y,z)\to(x,y,-z)$ and leaves every in-plane wavevector
unchanged, so it belongs to the little group of every $\bQ$ and commutes with
the Kohn-Sham, phonon and BSE Hamiltonians at every point of the zone. Every
non-degenerate state therefore carries an exact parity across the whole zone.
In a planar layer every atom lies in the mirror plane and maps onto itself; in
1H-WSe$_2$ the mirror maps W onto itself and swaps the two Se, $\pi=[0,2,1]$.

\subsection{Electrons}
For spinless states the diagonal of $D(\sigma_h;\bk)$ gives
$p_n(\bk)=\langle\psi_{n\bk}|\sigma_h|\psi_{n\bk}\rangle=\pm1$. With spin-orbit
coupling $\sigma_h=I\,C_{2z}$ acts on spin as $e^{-i\pi\sigma_z/2}=-i\sigma_z$, so
$\sigma_h^2=-1$ and the eigenvalues are $\pm i$. The label is
$\lambda_n=\mathrm{Im}\,p_n=\pm1$, chosen so that an exciton's parity is still a
product of labels (Sec.~\ref{sec:excparity}). WSe$_2$ gives
$\langle D(\sigma_h)^2\rangle=-1.0000$, confirming that the spin rotation is
included. At time-reversal-invariant points ($\Gamma$, $\mathbf{M}$) a spinor
band is a Kramers pair with opposite eigenvalues, so no single label exists
there. Time reversal maps $\bk\to-\bk$ and $+i\to-i$, so each band's label flips
between $\bk$ and $-\bk$ and every band is exactly 50\% even: at the
$\mathbf{K}$ band edges, where the orbitals are mirror-even $d$ states, the label
is the out-of-plane spin, and this is spin-valley locking written as mirror
parity.

\subsection{Phonons}
$\sigma_h$ acts on a displacement pattern $\mathbf{e}=(\mathbf{e}_\kappa)$ as
$(\sigma_h\mathbf{e})_\kappa=M\,\mathbf{e}_{\pi(\kappa)}$ with $M=\mathrm{diag}(1,1,-1)$
and $\pi$ the atom permutation (an involution). The parity of a normalised mode
is
\begin{equation}
p_{\nu\bq}=\langle\mathbf{e}_{\nu\bq}|\sigma_h\mathbf{e}_{\nu\bq}\rangle
=\sum_\kappa\mathbf{e}_{\nu\bq,\kappa}^{\dagger}\,M\,\mathbf{e}_{\nu\bq,\pi(\kappa)}=\pm1 .
\label{eq:phparity}
\end{equation}
For a planar layer, $\pi=\mathrm{id}$ and $p=1-2w_z$ with
$w_z=\sum_\kappa|e_{\kappa z}|^2$: out-of-plane modes are odd. In WSe$_2$ this
shortcut fails for 34.4\% of $(\bq,\nu)$: the Se breathing mode A$'_1$ moves
along $z$ yet is even, and the Se antiphase mode E$''$ moves in plane yet is odd.
Equation~(\ref{eq:phparity}) gives $\pm1$ to $7\times10^{-16}$ there. The
eigenvectors are taken from the electron-phonon database itself, in the order of
the couplings: branches labelled from a separate \texttt{matdyn.x} run,
ordered by its own frequencies, disagreed at 5\% of the points in GaN.

\subsection{Excitons}
\label{sec:excparity}
Applying a rotation $R$ to exciton $T$ and projecting on $S$,
\begin{equation}
D_{ST}(R)=\sum_{\bk}\sum_{vc,v'c'}\bigl[A^{S}_{vc\bk}\bigr]^{*}D^{c}_{cc'}(R;\bk)\,
A^{T}_{v'c',R^{-1}\bk}\bigl[D^{v}_{vv'}(R;\bk-\bQ)\bigr]^{*}.
\label{eq:Dexc}
\end{equation}
For $\sigma_h$, $R^{-1}\bk=\bk$, and for non-degenerate bands
\begin{equation}
\chi_S\equiv D_{SS}(\sigma_h)=\sum_{vc\bk}|A^S_{vc\bk}|^2\,p_c(\bk)\,p_v^{*}(\bk-\bQ),
\label{eq:chi}
\end{equation}
real and equal to $\pm1$ for spinless and spinor states alike (for spinors
$p_cp_v^*=(\pm i)(\mp i)$ or $(\pm i)(\pm i)^*$). Using \texttt{Yambo}'s
storage, the valence matrix must be taken at the hole's wavevector, $\bk+\bQ$ in
table indices. Of the four ways of attaching the shift, only this one gives
$|\chi_S|=1$ for every state (100\% of states at all 61 irreducible $\bQ$ in GaN,
against at most 35\% for the others); the wrong pairing produced a fictitious
ring of fractional parity around $\Gamma$.

\subsection{Selection rule}
The phonon perturbation transforms like its mode,
$\sigma_h\,\partial V_\nu\,\sigma_h^{-1}=p_\nu\,\partial V_\nu$. Inserting
$\sigma_h^\dagger\sigma_h=1$ twice,
\begin{equation}
\G_{\beta\lambda\nu}=\langle\beta|\sigma_h^\dagger(\sigma_h\partial V_\nu\sigma_h^{-1})\sigma_h|\lambda\rangle
=p_\beta^{*}\,p_\nu\,p_\lambda\,\G_{\beta\lambda\nu},
\end{equation}
so $\G_{\beta\lambda\nu}\neq0$ only if $p_\beta p_\nu p_\lambda=+1$: an even
phonon connects excitons of equal parity, an odd phonon excitons of opposite
parity.

\emph{Quantifying it} (\texttt{xphd selection-rule}). Each term
$|\G_{\beta\lambda\nu}(\bq)|^2$ is classed as allowed, forbidden or undefined
(a final state with $|\chi_\beta|$ below a tolerance). Results:
\begin{itemize}[nosep,leftmargin=1.2em]
\item GaN, bright state at $\Gamma$: forbidden/allowed $=5\times10^{-7}$.
\item WSe$_2$, dark (odd) and bright (even) states at $\Gamma$: $1.4\times10^{-4}$
and $1.0\times10^{-3}$ at tolerance 0.5, carried entirely by one pair of final
states 0.13~meV apart and of opposite parity, at one star of 12 $\bq$, returned
mixed by the solver ($|\chi|=0.915$, $0.951$). At tolerance 0.96:
$1.8\times10^{-7}$ and $1.4\times10^{-7}$, with 22\% and 49\% of the coupling then
undefined. With manifold rotation (Sec.~\ref{sec:manifold}), state~1:
$6.3\times10^{-6}$ with 5.9\% undefined.
\end{itemize}
\texttt{where\_forbidden.py} separates three fingerprints of residual forbidden
weight: a phonon swap (labels and couplings in different orders), a state mix (a
near-degenerate opposite-parity partner, or $|\chi|<0.95$), and a diffuse
admixture $\varepsilon=(1-|\chi|)/2$ (forbidden/allowed into a state scales as
$\varepsilon$, state by state).

\subsection{Irreducible representations at $\mathbf{K}$ and the emission rule}
The little group at $\bK$ is C$_{3h}$, Abelian, with six one-dimensional
irreducible representations (Table~\ref{tab:c3h}). Time reversal maps $\bK$ to
$\bK'$ and adds no degeneracy at $\bK$; $\sigma_vT$ fixes $\bK$ but maps each
C$_3$ eigenvalue to itself and squares to $+1$, so it adds none either. Which of
A, E$_1$, E$_2$ labels a state depends on the C$_3$ axis (through Ga, N or the
hexagon centre), because a Bloch sum at $\bK$ acquires a phase from sites off
the axis; the $\sigma_h$ parity does not.

\begin{table}[h]
\centering\small
\caption{Character table of C$_{3h}$, $\omega=e^{2\pi i/3}$.}
\label{tab:c3h}
\begin{tabular}{lcccccc}
\toprule
 & $E$ & $C_3$ & $C_3^2$ & $\sigma_h$ & $S_3$ & $S_3^5$\\
\midrule
A$'$ & 1 & 1 & 1 & 1 & 1 & 1\\
A$''$ & 1 & 1 & 1 & $-1$ & $-1$ & $-1$\\
E$'_1$ & 1 & $\omega$ & $\omega^2$ & 1 & $\omega$ & $\omega^2$\\
E$'_2$ & 1 & $\omega^2$ & $\omega$ & 1 & $\omega^2$ & $\omega$\\
E$''_1$ & 1 & $\omega$ & $\omega^2$ & $-1$ & $-\omega$ & $-\omega^2$\\
E$''_2$ & 1 & $\omega^2$ & $\omega$ & $-1$ & $-\omega^2$ & $-\omega$\\
\bottomrule
\end{tabular}
\end{table}

For the dark A$''$ exciton at $\bK$ to emit, a phonon must carry away the
momentum $\bK$ while a photon is emitted. The exciton must be contained in the
product of the phonon and the photon,
\begin{equation}
\Gamma_{\mathrm{exc}}\subset\Gamma_\nu\otimes\Gamma_{\mathrm{photon}},
\label{eq:emission}
\end{equation}
a form invariant under a change of the C$_3$ axis (the conjugated form
$\Gamma_\nu\subset\Gamma^*_{\mathrm{exc}}\otimes\Gamma_{\mathrm{photon}}$
renames the helicities but selects the same branch). In-plane light is E$'$ in
D$_{3h}$ and splits at $\bK$ into E$'_1\oplus$E$'_2$, one component per helicity.
For A$''$, a photon of E$'_1$ symmetry requires an E$''_2$ phonon and one of
E$'_2$ symmetry an E$''_1$ phonon. Only E$''_1$ occurs among the six GaN branches
at $\bK$ (Table~\ref{tab:phK}), so one branch mediates the emission and its light
has a single helicity; time reversal gives $\bK'$ the opposite one.

\begin{table}[h]
\centering\small
\caption{GaN phonons at $\bK$ (C$_3$ about the \tbd{Ga/N} site) and whether each
can mediate the emission of the A$''$ exciton under Eq.~(\ref{eq:emission}).}
\label{tab:phK}
\begin{tabular}{cccccc}
\toprule
Branch & $\hbar\omega$ (meV) & $\omega$ (cm$^{-1}$) & Irrep & $\sigma_h$ & Mediates\\
\midrule
1 & 16.2 & 130.8 & E$''_1$ & odd & yes\\
2 & 17.9 & 145.1 & A$''$ & odd & no\\
3 & 28.7 & 231.4 & E$'_2$ & even & no\\
4 & 42.4 & 341.8 & A$'$ & even & no\\
5 & 93.7 & 756.3 & E$'_1$ & even & no\\
6 & 100.4 & 810.3 & E$'_2$ & even & no\\
\bottomrule
\end{tabular}
\end{table}

\emph{When is the sideband fully polarised?} Let the C$_3$ eigenvalues of the
two flexural modes at $\bK$ be $\phi_1,\phi_2$ (two of $1,\omega,\omega^2$) and
that of the dark exciton $\varepsilon$. By Eq.~(\ref{eq:emission}) a photon of
eigenvalue $\omega$ needs a phonon of eigenvalue $\varepsilon\omega^2$, one of
eigenvalue $\omega^2$ a phonon of $\varepsilon\omega$. Both helicities have a
partner only if $\{\varepsilon\omega,\varepsilon\omega^2\}=\{\phi_1,\phi_2\}$,
i.e.\ only if $\varepsilon$ is the third eigenvalue. If $\varepsilon$ equals
$\phi_1$ or $\phi_2$, one helicity is available and each valley's sideband is
fully circularly polarised. Shifting the axis multiplies $\varepsilon,\phi_1,
\phi_2$ by the same phase, so the criterion is axis-independent. In GaN the two
flexural modes at $\bK$ are sublattice-polarised (each moves one sublattice),
E$''_1$ at 16.2~meV and A$''$ at 17.9~meV; the A$''$ exciton ($\varepsilon=1$)
coincides with the A$''$ mode, so only E$''_1$ mediates.

\subsection{Two classes of planar D$_{3h}$ monolayers}
In a planar honeycomb binary each band edge is $\pi$-like ($p_z$, odd) or
$\sigma$-like ($s$, $p_{x,y}$, even), and for an exciton dominated by one
transition $\chi=p_cp_v$. The valence maximum is the anion $p_z$ state at $\bK$
(odd); the class is set by the conduction minimum (Table~\ref{tab:classes}).
\emph{Class~I}: both edges $\pi$-like at $\bK$; the lowest exciton is the bright,
even E$'$ state at $\Gamma$ and its sidebands are in-plane optical replicas
(E$'\subset$E$'\otimes$E$'$; the flexural branch is forbidden). \emph{Class~II}:
a $\sigma$-like conduction band at $\Gamma$ lies lowest; the lowest exciton is
momentum-dark and odd at $\bK$, on a band disconnected from the bright one, and
its sideband needs an odd phonon at $\bK$.

\begin{table}[h]
\centering\small
\caption{The two classes. Rows marked \tbd{} await the fixed-code runs.}
\label{tab:classes}
\begin{tabular}{llllll}
\toprule
 & VBM & CBM & lowest exciton & sideband phonon & $|P|$\\
\midrule
hBN (I) & $\bK$, N $p_z$, odd & $\bK$, B $p_z$, odd & $\Gamma$, even, E$'$ (5.475 eV) & TO/LO, even & \tbd{0}\\
SiC (I) & $\bK$, C $p_z$, odd & $\bK$, Si $p_z$, odd & $\Gamma$, even, E$'$ & \tbd{TO/LO} & \tbd{0}\\
GaN (II) & $\bK$, N $p_z$, odd & $\Gamma$, Ga $s$, even & $\bK$, odd, A$''$ (3.429 eV) & E$''_1$, odd, 16.2 meV & 1.000\\
\bottomrule
\end{tabular}
\end{table}

% =====================================================================
\section{Near-degenerate manifolds}
% =====================================================================
\label{sec:manifold}

\subsection{Why near-degenerate states mix}
Let two states of opposite parity be separated by $\Delta$, and let the computed
BSE Hamiltonian break $\sigma_h$ by a small $\delta$ between them (numerical
noise). First-order mixing gives an admixture amplitude $\delta/\Delta$, i.e.\
$|\chi|=1-2(\delta/\Delta)^2$. For $\Delta=0.13$~meV, $\delta\sim10~\mu$eV
already gives $|\chi|\approx0.99$, and the WSe$_2$ pair came back at 0.915 and
0.951. Per-state parity then either counts a forbidden residue or excludes the
states.

\subsection{Splitting by the manifold's parity eigenstates}
Group the states at each $\bQ$ into manifolds of energies closer than a
tolerance (\texttt{xphd parity --manifold-tol}, default 5~meV) and compute, with
Eq.~(\ref{eq:Dexc}), the matrix $M_{mn}=\langle m|\sigma_h|n\rangle$ within each.
For a closed manifold $M$ is unitary and Hermitian, with eigenvalues $\pm1$;
its eigenvectors $v_i$ are the parity eigenstates the solver could equally have
returned. The state $\partial V_\nu|\lambda\rangle$ projected on the manifold is
$|\psi\rangle=\sum_n G_n|n\rangle$ with $G_n=\langle n|\partial V_\nu|\lambda\rangle$,
the archive's amplitudes. Its weight on eigenvector $i$ is
\begin{equation}
W_i=|\langle v_i|\psi\rangle|^2=\Bigl|\sum_n (v_i)_n^{*}\,G_n\Bigr|^2,
\qquad \sum_iW_i=\sum_n|G_n|^2 ,
\end{equation}
and $W_i$ is allowed or forbidden according to the eigen-parity $\mu_i=\pm1$.
Nothing is excluded for being mixed; only eigenvectors with $|\mu_i|$ below the
tolerance remain undefined. The weights in each parity sector,
$\langle\psi|\tfrac12(1\pm M)|\psi\rangle$, do not depend on the basis within the
manifold.

\emph{Rotated states.} An archive's final state at $\bQ+\bq$ is the image of its
irreducible parent, Eq.~(\ref{eq:rotate}). For spatial $R$,
$R^{-1}\sigma_hR=\sigma_h$ on excitons, so the manifold matrix is the parent's;
for an operation containing time reversal (antiunitary) it is
$M^{*}$. (For spinors, $\sigma_h$ anticommutes with $C'_2$ and $\sigma_v$ in the
double group; the electron and hole each acquire the sign and it cancels for
the exciton, which is why per-state exciton parity is constant over a star.)
The time-reversal flag per full-zone point is the one \texttt{xphd generate}
uses: symmetry index $\ge n_{\mathrm{sym}}/(1+t_{\mathrm{rev}})$.

\emph{Result.} In WSe$_2$ (state 1) the rotation lowers forbidden/allowed from
$1.4\times10^{-4}$ to $6.3\times10^{-6}$ without excluding anything beyond the 5.9\%
already undefined. Only 70\% of the eigen-parities lie within 0.02 of $\pm1$.

\subsection{Closure diagnostics}
An eigen-parity short of $\pm1$ means the manifold is not closed: its partner is
missing.
\begin{itemize}[nosep,leftmargin=1.2em]
\item \emph{Cut by the state count} (\texttt{manifold\_closure.py}): compare the
eigen-parities of each manifold cut at the archive's state count with those of
the complete manifold from a parity run with more states. In WSe$_2$, 75.8\%
within 0.02 either way: the cut is not the cause.
\item \emph{Partner beyond the tolerance}: widening the manifolds from 5 to
20~meV changed nothing (75.3\% to 75.0\%).
\item \emph{Window not closed} (\texttt{window\_closure.py}): $\sigma_h$ maps a band
partly outside the BSE window when its mirror partner lies outside it. Then the
row norm $\sum_{j\in\mathrm{window}}|D_{ij}(\sigma_h;\bk)|^2<1$, and every exciton
built on it maps partly outside the BSE space. \tbd{Run on WSe$_2$.}
\end{itemize}

% =====================================================================
\section{Exciton-phonon linewidths}
% =====================================================================
\label{sec:lw}

\subsection{Golden rule}
To first order in the coupling and in the dilute limit, the full width at half
maximum of state $\lambda$ at $\bQ$ is
\begin{equation}
\begin{split}
\Gamma_{\lambda\bQ}(T)=2\pi\sum_{\beta\nu}\int_{\mathrm{BZ}}\frac{A_{\mathrm{uc}}\,d^2q}{(2\pi)^2}\,
|\G_{\beta\lambda\nu}(\bQ,\bq)|^2
&\Bigl[(N_{\nu\bq}+1)\,\delta\bigl(E_{\lambda\bQ}-E_{\beta\bQ+\bq}-\hbar\omega_{\nu\bq}\bigr)\\
&+N_{\nu\bq}\,\delta\bigl(E_{\lambda\bQ}-E_{\beta\bQ+\bq}+\hbar\omega_{\nu\bq}\bigr)\Bigr],
\end{split}
\label{eq:fgr}
\end{equation}
the imaginary part of the Fan-Migdal self-energy, with the final-state exciton
occupation set to zero. Since the zone has area $(2\pi)^2/A_{\mathrm{uc}}$, the
measure is $A_{\mathrm{uc}}\,d^2q/(2\pi)^2=d^2\tilde q/\tilde A_{\mathrm{BZ}}=
(2/\sqrt3)\,d^2\tilde q$ for $\tilde q$ in units of $\bb$, and
$(1/N_q)\sum_{\bq}$ on a mesh. States within 2~meV are one manifold, their
widths averaged; the 2.8~meV exchange-split bright doublet of GaN is reported
state by state.

\subsection{Linear-triangle integration}
On the $24\times24$ mesh a Lorentzian or Gaussian $\delta$ would need a width
large enough to make the result depend on it, and in two dimensions it blurs the
kinematic thresholds. The $\delta$ functions are integrated analytically
\cite{Jepsen1971,Lehmann1972}: each mesh cell is split into two triangles, and in
each the energy difference $f=E_\lambda-E_\beta\mp\hbar\omega$ is linear between
its vertex values $f_1\le f_2\le f_3$. The values of a linear function over a
triangle of area $A_T$ are distributed with density
\begin{equation}
\rho_T(\varepsilon)=\frac{1}{A_T}\int_T\delta(\varepsilon-f)\,d^2r=
\begin{cases}
\dfrac{2(\varepsilon-f_1)}{(f_2-f_1)(f_3-f_1)}, & f_1\le\varepsilon<f_2,\\[1.5ex]
\dfrac{2(f_3-\varepsilon)}{(f_3-f_1)(f_3-f_2)}, & f_2\le\varepsilon<f_3,\\[1ex]
0 & \text{otherwise},
\end{cases}
\end{equation}
normalised to $\int\rho_T\,d\varepsilon=1$. The integral of Eq.~(\ref{eq:fgr}) is
then $\sum_T w_T A_T\rho_T(0)$, with $w_T$ the (piecewise-constant) weight
$|\G|^2(N+1)$ or $|\G|^2N$ of that triangle. This is done on a mesh refined
15-fold, to $360\times360$ (factor 30 for WSe$_2$), which needs the phonon and
exciton energies there.

\subsection{Inputs on the refined mesh}
\emph{Phonons.} Diagonalising the DFPT dynamical matrices directly with
\texttt{matdyn.x}, including the two-dimensional non-analytic correction that
removes the spurious LO-TO splitting at $\Gamma$ \cite{Sohier2017}. They reproduce
the source mesh to 1.7~$\mu$eV on all branches, apart from a 0.237~meV acoustic
residue at $\Gamma$ removed by the acoustic sum rule.

\emph{Exciton energies.} Exciton bands are indexed by ascending energy, so where
two bands cross each indexed band has a kink. A global Fourier interpolant
\cite{Koelling1986,Pickett1988} rings from every kink across the zone; at the
$\bK$ minimum of GaN it made the band 34\% too steep and gave it a spurious
twofold anisotropy. Each refined energy is instead taken from a quadratic about
the nearest BSE point,
\begin{equation}
E(u,v)=E_0+g_1u+g_2v+Au^2+Buv+Dv^2,
\end{equation}
with $E_0$ the BSE energy and $(u,v)$ the reduced offset. The six neighbours lie
at $\pm(h,0)$, $\pm(0,h)$, $\pm(h,-h)$. Sums over opposite pairs cancel the odd
terms and fix the curvature exactly:
\begin{equation}
S_1\equiv E(h,0)+E(-h,0)-2E_0=2Ah^2,\quad S_2=2Dh^2,\quad S_3=2(A-B+D)h^2,
\end{equation}
so $A=S_1/2h^2$, $D=S_2/2h^2$, $B=(S_1+S_2-S_3)/2h^2$. Differences give three
estimates of the gradient, $d_1=E(h,0)-E(-h,0)=2hg_1$, $d_2=2hg_2$ and
$d_3=E(h,-h)-E(-h,h)=2h(g_1-g_2)$, each up to third-order terms. Minimising
$(d_1-x)^2+(d_2-y)^2+(d_3-x+y)^2$ over $x=2hg_1$, $y=2hg_2$ gives
\begin{equation}
g_1=\frac{2d_1+d_2+d_3}{6h},\qquad g_2=\frac{d_1+2d_2-d_3}{6h}.
\end{equation}
The residual is the odd, trigonal-warping part of the band, third order in the
offset. The quadratic reproduces $E_0$ exactly and its neighbours to third
order, and because it uses only nearest neighbours a crossing affects only the
cells next to it. Quadratic interpolation onto a refined mesh followed by
linear-triangle integration is the hybrid tetrahedron scheme of MacDonald
\emph{et al.} \cite{MacDonald1979,Dong2025}, here in two dimensions
\cite{Wiesenekker1988}. We centre each quadratic on a mesh point rather than
fitting one per triangle through its vertices and edge midpoints: such
per-triangle quadratics are continuous but kinked at the vertices, and a band
minimum always lies on a vertex.

\emph{Tests.} On model bands with closed-form linewidths, the per-triangle
scheme erred by up to 17\% at a minimum and the present one by at most 2.4\%; a
crossing 150~meV above the minimum shifted the Fourier result by 35\% and the
present one by 2\%. At the GaN $\bK$ minimum, fitting the curvature to the
first, second and third shells of BSE points gives 12.44, 12.26 and
12.18~eV$\,\bb^{-2}$. Predicting every other BSE point from the rest gives an rms
error of \tbd{X}~meV for this scheme and \tbd{Y}~meV for Fourier interpolation
(\texttt{compare\_band\_interp.py}). A Wannier representation of the exciton
\cite{Haber2023} would treat crossings rigorously.

\emph{Couplings.} $|\G|^2$ is constant within each coarse cell, except for the
three acoustic branches in the cell around $\bq=0$, where the coupling vanishes
at the cell's own mesh point. There we use the long-wavelength form
$|\G|^2\,2\hbar\omega=Cq^2$ (i.e.\ $|\G|^2\propto q$ for a linear branch), with
$C$ fitted for each pair of initial and final states to the innermost shell of
the source mesh. Phonons below 0.5~meV are excluded. In \texttt{xphd linewidth}
this is \texttt{--acoustic-model q2 --n-shell 1}; the default is \texttt{none}.

\emph{Symmetry.} Each state's width is invariant under the point group, so
Eq.~(\ref{eq:fgr}) is evaluated at the irreducible $\bQ$ (61 for $24\times24$,
D$_{3h}$ with time reversal) and unfolded.

\subsection{The band minimum in closed form}
At the global minimum (GaN: $\bK$) no lower state exists; the width comes only
from absorption of acoustic phonons within the valley, into final states inside
the first mesh cell, where band and coupling are extrapolated. Model the valley
as $E(q)=E_{\min}+bq^2$, the branch as $\hbar\omega=\hbar cq$ and the coupling
as $|\G|^2\,2\hbar\omega=Cq^2$, with $q$ in units of $\bb$, following the
deformation-potential treatment of exciton-acoustic-phonon scattering
\cite{Kaasbjerg2013,Rudin1990,Selig2016}. Absorption from the bottom conserves
energy when $E_{\min}+\hbar cq=E_{\min}+bq^2$, i.e.\ at
\begin{equation}
q^*=\frac{\hbar c}{b},\qquad \hbar\omega^*=\hbar cq^*=\frac{(\hbar c)^2}{b}
\quad\bigl(=2Mc^2\ \text{for}\ b=\hbar^2/2M\bigr).
\end{equation}
Then, with the measure of Sec.~\ref{sec:lw}.1 and $d^2q=2\pi q\,dq$,
\begin{align}
\Gamma&=2\pi\cdot\frac{2}{\sqrt3}\int 2\pi q\,dq\;\frac{Cq}{2\hbar c}\,N(\hbar cq)\,
\delta\bigl(bq^2-\hbar cq\bigr)
=\frac{8\pi^2}{\sqrt3}\,\frac{Cq^{*2}}{2\hbar c}\,\frac{N(\hbar\omega^*)}{|2bq^*-\hbar c|}\nonumber\\
&=\frac{4\pi^2}{\sqrt3}\,\frac{C\,q^{*2}}{(\hbar c)^2}\,N(\hbar\omega^*)
=\frac{4\pi^2}{\sqrt3}\,\frac{C\,N(\hbar\omega^*)}{b^2},
\label{eq:closed}
\end{align}
using $|2bq^*-\hbar c|=\hbar c$. Above $\hbar\omega^*/k_B$,
$N\approx k_BT/\hbar\omega^*$ and $\Gamma\approx(4\pi^2/\sqrt3)\,Ck_BT/[b(\hbar c)^2]$
grows linearly in $T$; below it the width freezes out. For GaN all inputs come
from our data: $b=12.44$~eV$\,\bb^{-2}$, $c_{\mathrm{LA}}\approx\tbd{9.1}$~km/s,
$\hbar\omega^*=1.49$~meV ($q^*=0.011\,\bb$), and $C$ from the minimum-to-minimum
coupling (implying a deformation potential $\Xi\approx\tbd{8.4}$~eV; TA gives
$10^{-3}$ of LA and ZA nothing, by parity). Equation~(\ref{eq:closed}) gives
0.09, 10.6--11.4 and 45--48~meV at 5, 77 and 300~K (the range from extrapolating
$C$ to $q\to0$); the code gives 0.087, 10.4 and 50.5~meV, the excess at 300~K
being about 6~meV of absorption beyond the first cell. At 77~K the width exceeds
$\hbar\omega^*$: the scattering is quasi-elastic and the line not strictly
Lorentzian.

\emph{Leaving the state versus leaving the valley.} The absorbed phonon moves the
exciton by $|\bq^*|=0.011\,\bb$ inside its valley, with almost no change of
energy. Excluding the first cell (\texttt{--acoustic-model none}) leaves the rate
of scattering beyond it, the escape from the valley bottom: 0.044~meV at 77~K
($\hbar/\Gamma\approx15$~ps) and \tbd{X}~meV at 300~K.

\subsection{Off-diagonal self-energy and validity}
The off-diagonal elements of the exciton-phonon self-energy \cite{Chan2023}
change the absorption lineshape by less than \tbd{0.3}\%: the diagonal width of
the lowest bright exciton (51.9~meV at 77~K) exceeds the largest off-diagonal
element (\tbd{2.4}~meV) by a factor \tbd{22}. For the highest states, widths reach
several hundred meV (lifetimes $\sim1$~fs, comparable to the band spacing); there
Eq.~(\ref{eq:fgr}) no longer describes a sharp resonance.

\emph{GaN values.} Bright state at $\Gamma$: 36.7, 51.9 and 150.9~meV at 5, 77 and
300~K (12.7~fs at 77~K, 4.4~fs at 300~K; 99.8\% emission at 5~K, 43\% absorption
at 300~K); its partner, 56.8~meV at 77~K. $\bK$ minimum: 0.087, 10.4 (63~fs) and
50.5~meV (13~fs).

% =====================================================================
\section{Phonon-assisted luminescence and helicity}
% =====================================================================
\label{sec:pl}

\subsection{Second-order emission}
An exciton $\beta$ at momentum $\bq$ emits through a virtual bright state
$\lambda$ at $\bQ=0$, emitting or absorbing a phonon $(\mu,\bq)$ that carries the
momentum \cite{Cannuccia2019,Paleari2019}:
\begin{equation}
\begin{split}
I^{\mathrm{PL}}(\omega;T)\propto\frac{1}{N_q}\sum_{\beta,\bq,\mu,\pm}
&\Biggl|\sum_{\lambda}\frac{T_\lambda\,\G_{\lambda\beta\mu}(\bq,-\bq)}
{E_{\beta\bq}\mp\hbar\omega_{\mu\bq}-E_{\lambda\mathbf{0}}}\Biggr|^2\\
&\times N_{\beta\bq}(T_{\mathrm{exc}})\,F^{\pm}_{\mu\bq}(T)\,
\delta\bigl(\hbar\omega-E_{\beta\bq}\pm\hbar\omega_{\mu\bq}\bigr),
\end{split}
\label{eq:pl}
\end{equation}
with $T_\lambda$ the dipole of the bright state,
$N_{\beta\bq}\propto e^{-E_{\beta\bq}/k_BT_{\mathrm{exc}}}$ and
$F^\pm=N_{\mu\bq}+\tfrac12\pm\tfrac12$. In GaN ($T=T_{\mathrm{exc}}=10$~K, 15
states, six branches, 5~meV broadening) the spectrum below the direct line is a
single satellite at 3.4129~eV, 16.1~meV below the dark $\bK$ exciton
(3.4290~eV), matching the E$''_1$ branch at $\bK$ (16.2~meV); branch-resolved, it
exceeds every other branch by seven orders of magnitude.

\subsection{Helicity and valley}
Replace $T_\lambda$ by $\hat{\boldsymbol\epsilon}_\pm^{*}\cdot\mathbf{d}_\lambda$ with
$\hat{\boldsymbol\epsilon}_\pm=(\hat{\mathbf{x}}\pm i\hat{\mathbf{y}})/\sqrt2$ and the
complex Cartesian dipoles $\mathbf{d}_\lambda$ of the $\bQ=0$ excitons, and
separate the emitters $\beta$ by valley. The degree of circular polarisation is
$P=(I_+-I_-)/(I_++I_-)$. Because $\hat{\boldsymbol\epsilon}_\pm$ form an
orthonormal basis of the plane, for any amplitude vector $\mathbf{v}$,
\begin{equation}
|\hat{\boldsymbol\epsilon}_+^{*}\cdot\mathbf{v}|^2+|\hat{\boldsymbol\epsilon}_-^{*}\cdot\mathbf{v}|^2
=|v_x|^2+|v_y|^2 ,
\end{equation}
so $I_++I_-$ equals the helicity-summed spectrum exactly (a sum rule checked to
$3.7\times10^{-16}$ in GaN). Results: $P=-1.0000$ from $\bK$, $+1.0000$ from $\bK'$,
equal valley intensities to 0.12\% (ratio 0.9988). The dipoles must come from a
single BSE solution: with dipoles and eigenvectors from different runs, or with
valleys summed, the helicity is lost. In hBN (class I) no valley-locked
satellite is expected \tbd{(negative control to run)}.

\subsection{Exciton dipoles and oscillator strengths}
\texttt{xphd dipoles} builds $\mathbf{d}_\lambda=\sum_{vc\bk}A^{\lambda}_{vc\bk}\,
\mathbf{d}_{cv}(\bk)$ from the single-particle dipoles of \texttt{ndb.dipoles}
in the BSE window, rotated to the full zone like the envelopes (complex
conjugated for time-reversed points). It is validated against
\texttt{Yambo}'s oscillator strengths manifold by manifold: the ratio is the same
constant (the unit conversion) for every manifold, with exchange-split pairs
compared as sums. For the transport, \texttt{add\_osc.py} writes the in-plane
strengths $|d_x|^2+|d_y|^2$ into the $\Gamma$ archive (field \texttt{osc\_Q0});
$\sigma_h$-odd, $z$-polarised excitons receive none.

% =====================================================================
\section{Real-time exciton Boltzmann transport}
% =====================================================================
\label{sec:bte}

\subsection{Equation and rates}
For the occupations $F_i(t)$, $i\equiv(\lambda,\bQ)$
\cite{Zhou2021Perturbo,Chen2020},
\begin{equation}
\frac{dF_i}{dt}=-\sum_j\Bigl[P_{ij}F_i(1+F_j)-P_{ji}F_j(1+F_i)\Bigr],
\label{eq:bte}
\end{equation}
\begin{equation}
P_{ij}=\frac{2\pi}{\hbar}\frac{1}{N_q}\sum_\nu|\G_{\beta\lambda\nu}(\bQ,\bq)|^2
\Bigl[(N_{\nu\bq}+1)\,\delta(E_i-E_j-\hbar\omega_{\nu\bq})+N_{\nu\bq}\,\delta(E_i-E_j+\hbar\omega_{\nu\bq})\Bigr],
\end{equation}
with $\bq=\bQ_j-\bQ_i$. With the matrix elements, phonons and triangle weights of
Sec.~\ref{sec:lw}, the row sums are the linewidths, $\sum_jP_{ij}=\Gamma_i/\hbar$;
they agree to better than 0.1\% with Eq.~(\ref{eq:fgr}) evaluated on the same mesh
(without the in-cell acoustic channel, which moves an exciton within its own
cell and cannot move population between mesh states). The rate matrix couples
different momenta, so it is assembled on the full zone: $576\times5=2880$ states
for GaN.

\subsection{Detailed balance}
A transition $i\to j$ by emission is integrated over the triangle containing
$\bq$, the reverse $j\to i$ by absorption over the one containing $-\bq$. The
two energy-conserving surfaces are identical by time reversal, but linear
interpolation gives them different areas $A^{\mathrm{em}}_{ij}\neq A^{\mathrm{abs}}_{ji}$,
and the system then accumulates particles and never equilibrates. With the
symmetrised weight $B_{ij}=\tfrac12(A^{\mathrm{em}}_{ij}+A^{\mathrm{abs}}_{ji})$,
\begin{equation}
P_{ij}=B_{ij}|\G|^2(N+1),\quad P_{ji}=B_{ij}|\G|^2N
\ \Rightarrow\ \frac{P_{ij}}{P_{ji}}=\frac{N+1}{N}=e^{\hbar\omega/k_BT}=e^{(E_i-E_j)/k_BT}
\end{equation}
exactly, for every energy-conserving pair. The rates are written without
dividing by $N$, which would amplify weakly populated channels. Particle number
is conserved to $10^{-9}$.

\emph{Stationary state.} With
$F=1/(e^{(E-\mu)/k_BT}-1)$, $F/(1+F)=e^{-(E-\mu)/k_BT}$, so
$P_{ij}F_i(1+F_j)=P_{ji}F_j(1+F_i)$ is exactly the ratio above: every pair is
balanced, and the Bose-Einstein distribution is stationary (verified to
$10^{-16}$); at these densities it is Boltzmann.

\subsection{Refining the rates}
The full matrix cannot be refined like a linewidth (each refined $\bq$ is a new
state; the size grows as the fourth power of the refinement), but its row sums
are ordinary zone integrals. With $s_i$ the ratio of the row sum on a fivefold
finer mesh to the coarse one, the matrix is rescaled as
$P_{ij}\to P_{ij}\sqrt{s_is_j}$, which leaves $P_{ij}/P_{ji}$ unchanged. A
one-sided rescaling $P'_{ij}=s_iP_{ij}$ would not: detailed balance
$P'_{ij}F_i=P'_{ji}F_j$ then gives $F_i/F_j=(s_j/s_i)\,e^{-(E_i-E_j)/k_BT}$, i.e.\
$F_i\propto e^{-E_i/k_BT}/s_i$, the wrong equilibrium. Without the correction the
coarse rates are low by \tbd{2.5}\%; with it they match the refined reference to
\tbd{X}\%. The thermalisation time is set by the refined row sums; the
branching ratios vary smoothly with $\bq$.

\subsection{Initial condition, integration, diagnostics}
Optical injection fills $\bQ=0$:
$F_\lambda(0)\propto|\mu_\lambda|^2\exp[-(E_\lambda-E_{\mathrm{pump}})^2/2\sigma^2]$
\tbd{(pump as in the main text)}, normalised to the injected density; results are
unchanged from $10^{-4}$ to $10^{-2}$ per state (linear regime). Equation~(\ref{eq:bte})
is stiff and is integrated with an implicit BDF solver (relative and absolute
tolerances $10^{-7}$ and $10^{-12}$); occupations far above the minimum can dip to
$-10^{-15}$, an integrator artefact. Two decompositions are reported: the filling
rate $\sum_jP^{(c)}_{ji}F_j$ of a state by channel $c$, and the aggregate flux
$\sum_{ij}F_iP^{(c)}_{ij}$; a channel can fill many sparsely populated states yet
carry little flux.

\subsection{Results and the finite-temperature equilibrium}
GaN at 77~K: the population relaxes to 4.26~meV above the minimum, against a
Boltzmann value of 4.09~meV on the same states. Both lie below $k_BT=6.6$~meV, the
value for a two-dimensional continuum, because the nearest mesh states around
$\bK$ lie $b h^2=22$~meV above it. Of the population, 35, 50, 90 and 99\% reaches
the $\bK$, $\bK'$ valleys (within $0.12\,\bb$) by 30, 41, 98 and 172~fs; it then
escapes the valley bottom at 0.044~meV (15~ps), although each $\bK$ state loses
coherence in 63~fs. The valleys fill mainly through the fourth branch in its
in-plane, LA-like form: valley and feeding states are both $\sigma_h$-odd, so the
phonon must be even.

At a temperature where $k_BT$ is comparable to the energy of the first shell of
mesh points around the minimum, the equilibrium on the mesh is not concentrated
on the minimum's own point: a light exciton at 300~K can hold only
$\sim20\%$ there (\texttt{check\_bte.py} prints the share at the minimum's point and
in each valley against Boltzmann on the same states). A run has settled when it
matches Boltzmann; equal emission and absorption flux for every branch is the
same statement. The hBN run at 300~K behaves this way.

% =====================================================================
\section{Quasiparticle corrections: decomposition and use}
% =====================================================================
\label{sec:gw}

\subsection{Newton solution and its parts}
\texttt{Yambo}'s linearised quasiparticle equation is
\begin{equation}
E=E_o+Z\bigl[\Sigma_x+\mathrm{Re}\,\Sigma_c(E_o)-V_{xc}\bigr],\qquad
Z=\bigl[1-\partial_\omega\mathrm{Re}\,\Sigma_c|_{E_o}\bigr]^{-1},
\end{equation}
so a correction splits into an exchange part $Z(\Sigma_x-V_{xc})$ and a
correlation part $Z\,\Sigma_c(E_o)$. With \texttt{ExtendOut} every term is
printed (columns \texttt{Vxc}, \texttt{Vnlxc}$=\Sigma_x$, \texttt{Sc|Eo},
\texttt{Re(Z)}), and the equation reproduces the \texttt{E-Eo} column exactly.
Columns must be read by their header names: the layout is
\texttt{K-point, Band, Eo, E, E-Eo, \ldots}, and taking the fourth column for
$E-E_o$ once produced a false diagnosis (see the errata).
\texttt{qp\_decompose.py} and \texttt{cb\_splitting.py} read by name.

\subsection{WSe$_2$ without RIM-W}
Band edges (\texttt{cb\_splitting.py}): the gap is \emph{direct} at $\bK$, 2.422~eV
(VBM band 62 at $-1.1195$~eV, CBM band 64 at 1.3021~eV); the Q valley (the
$\Gamma$-$\bK$ midpoint) lies 0.451~eV above $\bK$ (DFT: 0.064~eV), so the
indirect K$\to$Q gap is 2.872~eV. The decomposition, in eV:
\begin{center}\small
\begin{tabular}{lccc}
\toprule
difference & total & exchange $Z(\Sigma_x-V_{xc})$ & correlation $Z\Sigma_c$\\
\midrule
gap at $\bK$ (62$\to$63) & $+1.181$ & $+3.135$ & $-1.954$\\
Q against $\bK$ (band 63) & $+0.368$ & $+0.178$ & $+0.191$\\
band 64 $-$ band 63 at $\bK$ & $-0.055$ & $-0.196$ & $+0.140$\\
\bottomrule
\end{tabular}
\end{center}
The gap behaves normally: correlation closes about 60\% of the exchange opening.
Two features disagree with the established picture of WSe$_2$. Both parts push
Q up against $\bK$, whereas the reference picture has Q at or below $\bK$ (an
indirect K$\to$Q gap, the lowest exciton the K-Q state at $\mathbf{M}$). And the
conduction spin splitting at $\bK$ reverses ($+37.2$~meV in DFT, $-18.2$~meV in
GW): exchange stabilises band 64, which has the spin of the valence maximum,
whereas the dark exciton of WSe$_2$ lying below the bright one requires the
opposite-spin band lower. With spin-orbit coupling, 200 bands in the response
and the self-energy amount to about 100 spinless bands; band convergence, and
the lattice constant (3.322~\AA\ against 3.28~\AA\ experimental), are the open
suspects \tbd{(tests pending)}.

\subsection{Using corrections in the BSE}
A rigid scissor shifts every conduction band equally and keeps DFT's valley
offsets and spin splittings: adequate within a valley, wrong between valleys.
\texttt{KfnQPdb} applies the $k$-dependent corrections correctly; the check is
the report's \texttt{[K+QP]} gap, which must equal the Kohn-Sham gap plus the
correction at that $k$. For WSe$_2$, RIM-W had to be off in both the GW and the
BSE (it removed the binding in the BSE kernel and reduced the GW gap correction
at $\bK$ from 1.18 to 0.10~eV); GaN and hBN ran correctly with it.

% =====================================================================
\section{WSe$_2$: deferred, and how to resume}
% =====================================================================
\label{sec:wse2}

WSe$_2$ is out of the GaN/hBN paper and kept in the package for a later paper
on transition-metal dichalcogenides. Paths: \texttt{D:\textbackslash Data\textbackslash WSe2\textbackslash}
with \texttt{LELPH} (SAVE, \texttt{ph.save}, \texttt{ndb.elph}), \texttt{GW},
\texttt{BSE}, \texttt{EXCPH} (Dmats, archives, parity files) and the phonon
folder with \texttt{copy.sh} and the \texttt{wse2\_matdyn} inputs.

\subsection{What is established}
\begin{itemize}[nosep,leftmargin=1.2em]
\item Spin-orbit symmetry machinery: $\langle D(\sigma_h)^2\rangle=-1.0000$; band
labels from $\pm i$, undetermined only at the Kramers points; phonon parity by
permutation ($\pi=[0,2,1]$), exact to $7\times10^{-16}$; every $\bQ=0$ exciton
$\pm1$.
\item Archive symmetry: $1.2\times10^{-3}$ (time reversal), $1.0\times10^{-3}$
(C$_3$).
\item Selection rule: $1.8\times10^{-7}$ and $1.4\times10^{-7}$ at tolerance
0.96; $6.3\times10^{-6}$ with manifold rotation (state 1). The residue at the
default tolerance is one solver-mixed pair; 25\% of manifold eigen-parities stay
short of $\pm1$ for a reason not yet found (state cut and manifold width ruled
out; window closure to test).
\item For WSe$_2$, RIM-W must be off in the GW and the BSE (it works for GaN and
hBN; the cause is not isolated); \texttt{KfnQPdb} works.
\end{itemize}

\subsection{What is open}
\begin{itemize}[nosep,leftmargin=1.2em]
\item The GW (no RIM-W, 200 bands, $a=3.322$~\AA, $12\times12$): direct gap at $\bK$
(2.42~eV), Q 0.43--0.45~eV above $\bK$, conduction spin splitting at $\bK$ reversed
($+37$~meV in DFT, $-18$~meV after GW).
\item Two published pictures. Mishra \emph{et al.} [PRB \textbf{98}, 045143 (2018)]:
fully relativistic ONCV with W 5$s$5$p$4$f$ semicore, PBE, $a=3.3175$~\AA, 18$\times$18
DFT mesh; G$_0$W$_0$ gap 2.19~eV, direct; Q lifted from 44 to 394~meV above $\bK$ (as
here); conduction splitting at $\bK$ kept at $+40$~meV. Chen, Sangalli and Bernardi
[PRR \textbf{4}, 043203 (2022)]: $a=3.27$~\AA; exciton minima at $\mathbf{M}$ and Q.
\item Reading: the Q--$\bK$ offset agrees with Mishra \emph{et al.}; the difference from
Bernardi is mainly the lattice constant (compressive strain lowers Q relative to $\bK$,
about 100~meV per \% in the bilayer), i.e.\ pseudopotential and functional. The
reversed conduction splitting is ours alone: an exchange effect ($Z\Delta\Sigma_x=-0.196$~eV
between the spin partners) to be checked against the W semicore and the $\Sigma_x$
cutoff. The gap (2.42 against 2.19~eV) is a mesh question.
\item Consequently the exciton landscape, the archives and everything downstream
are not yet trustworthy for WSe$_2$.
\end{itemize}

\subsection{Resume, in order}
\begin{enumerate}[nosep,leftmargin=1.4em]
\item \textbf{Pseudopotential and exchange} (cheap, settles the conduction order): check
the W valence configuration in the UPF header (5$s$5$p$ and 4$f$ semicore?); run
Hartree-Fock only (\texttt{yambo -x}) at $k=19$, bands 62--64, with \texttt{EXXRLvcs}
at 70, 100 and 140~Ry and watch $\Sigma_x(64)-\Sigma_x(63)$.
\item \textbf{Lattice constant}: DFT bands at $a=3.27$ and 3.28~\AA\ (scf + nscf only) to
see how far Q moves; G$_0$W$_0$ at Q and $\bK$ only if it changes the picture.
\item \textbf{GW band convergence} at $k=13$ (Q) and 19 ($\bK$) only:
\texttt{QPkrange} \texttt{13|13|62|64|} and \texttt{19|19|62|64|}; \texttt{BndsRnXp}
and \texttt{GbndRnge} at 200, 400, 600 (spinor bands; nscf with enough bands);
\texttt{rim\_cut} on, \texttt{rim\_w} off. After each: \texttt{cb\_splitting.py}
(Q$-$K offset; 64$-$63 at $\bK$) and \texttt{qp\_decompose.py}.
\item \textbf{Mesh}: G$_0$W$_0$ on $18\times18$ (Q and $\bK$ are mesh points) for the gap.
\item \textbf{Full GW} at the converged settings, all 19 $k$, bands 59--66.
\item \textbf{BSE} with \texttt{KfnQPdb= "E < GW/ndb.QP"}, no \texttt{rim\_w},
\texttt{BSSmod= "d"}, all 19 $\bQ$. Check the report's \texttt{[K+QP]} gap against
\texttt{cb\_splitting.py}; \texttt{exc\_transitions.py} at $\bQ=\mathbf{M}$
(\texttt{ndb.BS\_diago\_Q7}): is the lowest state a K-$\Lambda$ exciton?
\item \textbf{Symmetry}: \texttt{xphd parity --manifold-tol 0.005}; \texttt{window\_closure.py}
(\texttt{SIGMA\_H\_OP = 6}); \texttt{manifold\_closure.py}.
\item \textbf{Archives and analysis}: \texttt{xphd dmats}; \texttt{xphd generate}
(states enough to close the manifolds); \texttt{check-archive}; \texttt{labels\_from\_elph.py};
\texttt{selection-rule --rotate}; then linewidths, \texttt{add\_osc.py}, the BTE and
the maps as in Part~I, Sec.~9.2 (read with the errata).
\end{enumerate}

% =====================================================================
\section{Validation summary}
% =====================================================================
\label{sec:summary}

\begin{table}[h]
\centering\small
\begin{tabular}{p{6.2cm}p{4.6cm}p{5.4cm}}
\toprule
Check & Result & Where\\
\midrule
Exciton parity, $|\chi|=1$ & 100\% of states (GaN) & Sec.~\ref{sec:excparity}\\
Selection rule, forbidden/allowed & $5\times10^{-7}$ (GaN); $1.8\times10^{-7}$ (WSe$_2$, tol 0.96) & Sec.~\ref{sec:parity}\\
Star uniformity $T_{\max}/T_{\min}$ & 1.0001 (GaN) & Sec.~\ref{sec:excph}\\
Manifold-resolved symmetry, TR / C$_3$ & $1.3/3.3\times10^{-3}$ (GaN); $1.2/1.0\times10^{-3}$ (WSe$_2$) & Sec.~\ref{sec:excph}\\
Phonons, matdyn vs source mesh & 1.7~$\mu$eV & Sec.~\ref{sec:lw}\\
$\bK$ minimum vs closed form, 77~K & 10.4 vs 10.6--11.4~meV & Sec.~\ref{sec:lw}\\
Off-diagonal self-energy & $<\tbd{0.3}$\% of the lineshape & Sec.~\ref{sec:lw}\\
Helicity sum rule; $P$ per valley & $3.7\times10^{-16}$; $\mp1.0000$ & Sec.~\ref{sec:pl}\\
BTE row sums vs linewidths & $<0.1$\% & Sec.~\ref{sec:bte}\\
BTE detailed balance; stationarity & machine precision; $10^{-16}$ & Sec.~\ref{sec:bte}\\
BTE end state vs Boltzmann, 77~K & 4.26 vs 4.09~meV & Sec.~\ref{sec:bte}\\
Spinor mirror, $\langle D(\sigma_h)^2\rangle$ & $-1.0000$ (WSe$_2$) & Sec.~\ref{sec:parity}\\
Phonon parity by permutation & $\pm1$ to $7\times10^{-16}$ (WSe$_2$) & Sec.~\ref{sec:parity}\\
\bottomrule
\end{tabular}
\end{table}

\clearpage
\phantomsection
\addcontentsline{toc}{section}{References}
\bibliographystyle{unsrtnat}
\bibliography{partII}


\end{document}
